\section[Local unicyclic construction]{Local unicyclic discovery with two fixed cycle ports}
\label{sec:op3-local-unicyclic}

\paragraph{Status and model.}
The construction below is \statusproved{} as a proof draft awaiting
independent review. It extends \cref{thm:op3-local-top-tree} to arbitrary
finite simple connected unweighted unicyclic graphs, with an arbitrary
physical seed and arbitrarily large tree attachments. Neither the cycle
nor any attachment is supplied. The exact-real word and local graph-access
models are unchanged. The height-bounded top-tree balancing algorithm is
the same explicit \statussource{} import. The local state machine,
application algebra and named queries are implemented, while the audit
rebuilds reference hierarchies and counts those rebuilds separately.
The result does not settle arbitrary-graph OP3.

\begin{lemma}[Named coordinates through current cluster records; \statusproved]
\label{lem:op3-cluster-point-query}
Given a current height-$h$ cluster hierarchy, solved values at its at most
two root ports, and a stored incident-edge locator for an active vertex
$i$, its exact coordinate can be queried in $O(1+h)$ word operations.
The query uses a constant-size affine row and current parent pointers; it
does not recover or scan any other active coordinate. It therefore costs
$O(\log(2+n))$ on the height-bounded source hierarchy.
\end{lemma}
\begin{proof}
Start with the basis row selecting $i$ from its incident base edge's two
endpoints. A one-port base edge has one forget record, which immediately
pulls this row back to the structural leaf's port. Each structural join
introduces at most two application records. Compression recovers the shared
port as $w_m=\eta w_L+\zeta w_R+\sigma$; forgetting recovers
$w_b=\eta w_a+\sigma$; raking preserves each child port among the parent
ports. Hence a child's at most two port coordinates are affine functions
of the parent's at most two ports, computable through these at most two
new records in constant work. Pull the current affine row through that map
and follow the structural parent pointer. At the root, evaluate the row at
the solved port values and add $C$ if physical $u_i$ is requested.

The published source stores edge-leaf and structural parent pointers
(Section 2.4 of \cite{alstrup2005toptrees}). Store one current incident home
edge per vertex as in \cref{lem:op3-top-tree-payload}. An endpoint lookup,
including the root's first edge, therefore needs no search. Current child
to parent affine maps can be computed from the saved helper records or
cached with each freshly created parent in constant additional work.
Changed payload ancestors recompute these maps from current children.
No old numerical parent version is reactivated after an exposure update.
Equivalently, the expanded application hierarchy has at most twice the
structural height, and the query follows its current parent links.
The audit's whole-hierarchy index construction is a reference convenience;
it is not needed by this source-pointer implementation.
\end{proof}

\begin{lemma}[The unique exceptional boundary candidate; \statusproved]
\label{lem:op3-unicyclic-exception}
Let $U$ be a connected active set in a graph of cycle rank at most one.
If $G[U]$ is a tree, every inactive boundary vertex has one active neighbor,
except possibly one vertex having two active neighbors. If $G[U]$ is
already unicyclic, every inactive boundary vertex has exactly one active
neighbor. The exceptional vertex, if present, can be discovered from
cached incidences while scanning newly admitted rows.
\end{lemma}
\begin{proof}
Adjoin any collection $B$ of inactive boundary vertices and only their edges
to $U$. The resulting connected subgraph has cycle rank
\[
 |E(U)|-|U|+1+\sum_{j\in B}\bigl(|N(j)\cap U|-1\bigr)\leq1.
\]
Every summand is nonnegative. The assertions follow immediately, including
the impossibility of a third active neighbor. Since every active row has
already been scanned, a boundary record contains all its current active
parents. A newly scanned incidence either creates its first parent or
reveals the sole possible second parent. Thus no inactive row or unknown
cycle information is required.
\end{proof}

\paragraph{State before closure.}
Use the uniform shift of \cref{lem:op3-uniform-negative-shift} from the
start. Root the active spanning tree at the physical seed $v$, exposing
only $v$ at each checkpoint. Every ordinary candidate $j$ has one active
parent $i$, so its heap key is $\kappa_j=\lambda d_j/\gamma-C$.
Home only the minimum live key from each active vertex's heap. This gives
the root's exact ordinary sign certificate by
\cref{lem:op3-shifted-reporters}.

When a second active parent is discovered for $j$, remove its entry from
the first parent's ordinary heap and exclude it from the new parent's
heap. Refresh the first parent's home edge. Store the exceptional triple
$(j,p,q)$ separately. Lazy heap deletion is also valid: a live-entry test
checks the cached parent tuple and active status; every stale entry is
removed at most once. At every checkpoint, including unsuccessful checks,
use \cref{lem:op3-cluster-point-query} twice to evaluate
\[
 g_j=\gamma(u_p+u_q)-\lambda d_j.
\]
Either an ordinary reporter or this exceptional gate may supply the next
strict positive admission. Stop before closure only if both are quiet.
The exceptional record is created at most once: it persists until its
vertex is admitted, which closes the sole cycle, or until termination.

\paragraph{The single closure and subsequent growth.}
Admitting an ordinary candidate adds one leaf to the stored spanning tree.
Admitting the exceptional candidate $j$ adds two edges to the active graph.
Keep either new edge in the spanning tree and designate the other as the
extra edge $\{c,j\}$. Re-root and rebuild the cached active spanning tree
once at geometric anchor $c$, re-homing its current minimum rows. Insert
its edges through the source interface in parent-before-child order, then
expose $c,j$. The uniform negative-load shift allows the physical seed to
be interior, and \eqref{eq:op3-one-edge-restoration} gives the exact full
active solution at the two retained ports. The underlying uncorrected
spanning-tree response remains a signed affine object; no positivity or
obstacle truncation is imposed on that response.

By \cref{lem:op3-unicyclic-exception}, every subsequent inactive candidate
has one active parent. Its admission adds a leaf to the stored spanning
tree and changes its parent's minimum row. Expose the same two fixed cycle
ports after each link. The root hull now gives a complete sign certificate
for all remaining ordinary boundary candidates, using the corrected root
values. Link operations may temporarily reset external vertices: apply the
unavailable-anchor convention of \cref{lem:op3-source-interior-adapter}
and make no gate query until the required exposure has been restored.
The source splits all clusters whose boundary set changes before an update,
so both one-port and two-port exposures preserve the callback contract.

\begin{theorem}[Local exact obstacle and ACL solves on unicyclic graphs; \statusproved]
\label{thm:op3-local-unicyclic}
For a finite simple connected unweighted graph with at least two vertices
and cycle rank at most one, the preceding deterministic local algorithm
returns the exact obstacle solution for
$\boldsymbol M=\boldsymbol D-\gamma\boldsymbol A$ and
$\boldsymbol b=\boldsymbol e_v-\lambda\boldsymbol d$.
For $0<\alpha<1$, $\lambda>0$, and returned positive support $U$, its total
exact-real word work is
\[
 O\!\left(1+\operatorname{cvol}(U)
       \log^4(2+\operatorname{cvol}(U))\right).
\]
Its space, even retaining all historical application versions, is
$O(1+\operatorname{cvol}(U)\log^3(2+\operatorname{cvol}(U)))$ words.
For $\lambda=\epsappr/2$ and $0<\epsappr<1$, it returns
$\widehat{\boldsymbol\pi}=\bar\alpha\boldsymbol D\boldsymbol u$ with
original residual between zero and $\lambda\boldsymbol d$, in fully
charged local work $\widetilde O(1/\epsappr)$.
\end{theorem}
\begin{proof}
If $\lambda d_v\geq1$, return zero after the seed degree query. Otherwise
start at the positive singleton. Exact Schur states and reporters,
\cref{lem:op3-unicyclic-exception}, and the separately checked exceptional
gate ensure that every admission has a strictly positive original gate.
The usual positive inverse block update makes the new coordinate positive
and increases the old coordinates. Thus every checkpoint is a positive
face, all admissions lie in the final obstacle support, and termination
occurs after at most $|U|$ admissions. At termination every active equation
is exact and every inactive gate is nonpositive. These are the obstacle KKT
conditions, proving uniqueness and correctness. Final recovery visits only
the current application records and adds $C$ once per emitted coordinate.
The original residual is nonnegative, equals $\lambda d_i$ on $U$, and is
at most $\lambda d_i$ outside $U$.

Write $V=\operatorname{cvol}(U)$. Each admitted original row is scanned once,
for at most $V$ adjacency entries. Distinct degree replies, discovered-label
records, active-membership tests and cached boundary incidences require
$O(1+V)$ logical operations; comparison dictionaries add a logarithm.
There is one heap insertion per first boundary incidence and at most one
later deletion per entry, with $O(\log(2+V))$ work per heap operation.
Each admission changes at most two old home payloads and one new payload;
the exceptional extra old-parent refresh occurs only once. By
\cref{lem:op3-top-tree-payload}, these changes and the source's link/expose
callbacks cost $O((1+V)\log^4(2+V))$ work. At most two named-coordinate
queries per checkpoint add $O(|U|\log(2+|U|))$ work. Constant-size root
solves and the final failed ordinary and exceptional queries are included.

Closure re-roots only the cached active spanning tree, at most once.
Traversing its edges, assigning new depths and homes, and rebuilding through
the source interface costs $O(|U|\log^4(2+|U|))$, with
$O(|U|\log^3(2+|U|))$ newly allocated persistent state. This is charged to
the final $V$, even when the remaining graph is enormous. The capacity-
doubling alternative in \cref{lem:op3-top-tree-payload} adds the same bounds;
no array is initialized using the ambient vertex count. All historical
application records satisfy the stated total allocation bound, while each
current hierarchy has only linear structural state and constant-size helper
records. Temporary arithmetic matrices and query rows have constant size;
their work is charged and they need not be retained after use. The final
recovery and dictionary output cost $O((1+V)\log(2+V))$.

For nonempty $U$, the source-mass identity gives $\lambda\operatorname{vol}(U)<1$.
The positive-degree graph implies $|U|\leq\operatorname{vol}(U)$, so
$V=O(1/\lambda)$. Substituting $\lambda=\epsappr/2$ proves the ACL bound.
\end{proof}

\paragraph{RPPR consequence.}
With $\lambda=\rho$, the same construction gives the exact regularized
solution $\boldsymbol x=\bar\alpha\boldsymbol D^{1/2}\boldsymbol u$,
with objective gap zero and $\widetilde O(1/\rho)$ work in the nonzero
regime $0<\rho<1/d_v$. This removes both the cycle-seed assumption and the
polynomial attachment-size factor in the earlier bounded-attachment result.
It remains a structural exact-word theorem, not a numerical implementation
or an arbitrary-graph OP3 theorem.

\paragraph{Measured and implementation boundary.}
\texttt{cluster\_point\_queries.py} checks 58,816 exact named coordinates
on all core trees through six vertices, including arbitrary root assignments,
one-port base edges, all recovery helper shapes and retained payload versions.

\texttt{local\_unicyclic\_clusters.py} uses actual strict local admissions on
all 78 connected tree and unicyclic atlas graphs through seven vertices,
every seed, four parameter pairs and two admission/extra-edge policies.
Independent dense references check every face, every original boundary
gate, cycle correction, final obstacle solution and ACL residual. There are
3,888 atlas executions and sixteen larger explicit attachment executions,
checking 16,371 faces and 26,043 original gates. They include 319 negative
uncorrected spanning-tree faces. Four finite implicit private-star examples
have 9, 17, 15 and 16 active vertices, respectively; their original row
scans inspect only 27, 51, 45 and 48 entries despite enormous ambient size.
Inactive hub rows are never scanned. Exact totals and source hashes are
recorded in \texttt{LOCAL\_UNICYCLIC\_CLUSTER\_AUDIT.json}.
The audit's rebuilt hierarchies and full validation recoveries are explicitly
separate reference work. The complexity theorem imports the verified online
source hierarchy, rather than claiming that these reference rebuilds are
fast online updates.
